\documentclass[10pt,a4paper]{amsart}
\usepackage[margin=19mm]{geometry}
\usepackage{amsmath,amssymb,mathtools}
\usepackage[scr=rsfs,cal=euler]{mathalfa}
\usepackage{xfrac}
\usepackage[unicode,hidelinks,bookmarks=false]{hyperref}
\newcommand{\ie}{i.e.\ }
\newcommand{\cf}{cf.\ }
\newcommand{\resp}{respectively\ }
\newcommand{\Subsection}{\subsection}
\providecommand{\nobreakdash}{\nobreak\hbox{-}}
\setlength{\emergencystretch}{2em}
\multlinegap0pt
\usepackage{amssymb}
\usepackage{enumerate}
\usepackage{xcolor}
\usepackage{tikz}
\newtheorem{lemma}{Lemma}
\newtheorem{theorem}{Theorem}
\newtheorem{remark}{Remark}
\newcommand{\kc}{\mathcal{C}}
\title{How far are $d$-dimensional copulas with uniform $(d-1)$-marginals from (total) independence?}
\author{Damjana Kokol Bukov\v{s}ek, Nik Stopar, Wolfgang Trutschnig}
\date{Source: arXiv:2608.18286v1 (CC BY 4.0)}
\begin{document}
\maketitle

\noindent\emph{Source and licence.} How far are $d$-dimensional copulas with uniform $(d-1)$-marginals from (total) independence?. Version arXiv:2608.18286v1 (CC BY 4.0), distributed by arXiv under the Creative Commons Attribution 4.0 International licence, http://creativecommons.org/licenses/by/4.0/, as recorded in the arXiv licence metadata for this exact version and shown on the abstract page https://arxiv.org/abs/2608.18286v1. Full licence notice: \texttt{licenses/arxiv-2608.18286-CC-BY-4.0.txt}.\par\smallskip

\noindent\textit{Editorial scope.} Counted unit: the article's theorem thm:mainD1 (source0.tex lines 266--268). The article's complete original proof of it is delivered with it (source0.tex lines 269--301, one \texttt{proof} environment of 33 lines). No mathematics is written, repaired or re-derived: the statement and the proof are the article's own lines, in the article's own notation, and no retained line is substituted or re-flowed.  Retained uncounted as the source-local context the proof needs: lines 231--238; lines 317--330; lines 332--360.  One declared adaptation: exactly 2 ancillary strings inside the retained spans are omitted (image-inclusion commands and, where the source repeats a label, the repeated label definitions) and no other line differs from the source; the figure environments, with their placement, captions and labels, are kept, and the package records the omitted strings in fidelity receipts bound to the affected bodies.  No retained line contains a citation, so the wrapper carries no bibliography and no bibliography entry.  The retained lines contain the article's own diagram code, which the wrapper typesets with the standard tikz packages; no image file is delivered and the package declares no asset dependency.  Editorial formatting only, in addition: an AMS \texttt{amsart} preamble that restates the article's own declarations and macros used by the retained lines, namely the environments \texttt{lemma}, \texttt{theorem}, \texttt{remark} on separate counters, together with the standard packages those lines need.  The retained environments print in this document as Lemma 1, Theorem 1, Remark 1 in order of appearance, whereas the article prints the same items with the numbering of its own full document.  The complete native source of the article as distributed in the arXiv e-print for this exact version is bundled unchanged as \texttt{sources/207979/source0.tex}.
\begin{lemma}\label{lem:haar}
    Suppose that $X$ and $Y$ are independent random variables. 
    If $X$ is uniformly distributed on $[0,1]$,  
    then so is the random variable $Z$ given by 
    $$
    Z:= X +Y \,\, \textup{mod (1)}.
    $$
\end{lemma}

\begin{figure}[ht]
\begin{center}   
  
  
\end{center}
\caption{Left panel: Sample of size $n=5.000$ of the copula $A \in \kc_3^{\Pi_2}$ (in dimension $d=3$) as 
    considered in the proof of Theorem 
    \ref{thm:mainD1} in magenta, bivariate marginal samples in gray (i.e., the 
    three-dimensional magenta points are projected onto the hyperplanes $h_1: z=0$,
    $h_2: y=1$, $h_3: x=0$). \\
    Right panel: Sample of the symmetric copula $B\in \kc_3^{\Pi_2}$ considered in the second 
    part of Remark \ref{rem:symm}.}
    \label{fig1}
\end{figure}

\begin{remark}\label{rem:symm}
    Notice that fixing $a \in (0,1)$, slightly changing the definition 
    of $X_d$ according to eq. (\ref{eq:defX_q}) to 
    \begin{equation*}
        X_d:= (X_1 + X_2 + \ldots + X_{d-1} + a)  \,\, \textup{mod (1)}, 
    \end{equation*}
    and proceeding like in the proof of Theorem \ref{thm:mainD1}
    results in another copula $A_a \in \kc_d^{\Pi_{d-1}}$ 
    which is completely dependent and which also fulfills 
    $D_1(A_a,\Pi_d)=\frac{1}{3}=\Delta_{D_1}$. 
    In other words, we can easily construct uncountably many elements in 
    $\kc_d^{\Pi_{d-1}}$ having $D_1$-distance $\frac{1}{3}$ to $\Pi_d$.
    Moreover, our construction can easily be modified in order to obtain symmetric 
    (exchangeable) copulas with the afore-mentioned maximization property
    w.r.t.~$D_1$. 
    In fact, replacing (\ref{eq:defX_q}) by 
    \begin{equation*}
        X_d:= 1 - \big((X_1 + X_2 + \ldots + X_{d-1})  \,\, \textup{mod (1)}\big), 
    \end{equation*}
    we obviously have $(\sum_{i=1}^d X_i) \,\,\textup{mod (1)} =0$ almost surely and for every 
    $j \in [d]$ the identity 
    $$
     X_j = 1 - \left(\left(\sum_{i\neq j} X_i\right)  \,\, \textup{mod (1)}\right)
    $$
    holds almost surely. As a direct consequence, the copula $B$ of $\mathbf{X}=(X_1,\ldots,X_d)$ is
    exchangeable and fulfills $D_1(B,\Pi_d)=\frac{1}{3}=\Delta_{D_1}$. 
    The right panel of Figure \ref{fig1} depicts a sample (plus bivariate marginal samples) 
    of the copula $B$ for $d=3$.
\end{remark}

\begin{theorem}\label{thm:mainD1}
For every $d \geq 3$ we have $\Delta_{D_1}=\frac{1}{3}$ and the supremum is attained.
\end{theorem}

\begin{proof}
    Considering the afore-mentioned properties of $D_1$, it suffices to show that the class 
    $ \kc_d^{\Pi_{d-1}}$ contains completely 
    dependent copulas, which can be done as follows.
    Let $X_1,\ldots,X_{d-1}$ denote independent, uniformly $[0,1]$-distributed random variables
    and set 
    \begin{equation}\label{eq:defX_q}
        X_d:= (X_1 + X_2 + \ldots + X_{d-1})  \,\, \textup{mod (1)}. 
    \end{equation}
    Then, as a direct consequence of Lemma \ref{lem:haar}, $X_d$ is uniformly distributed on 
    $[0,1]$ too. Letting $A$ denote the copula
    of the vector $\mathbf{X}=(X_1,\ldots,X_d)$ we obviously have that 
    $A \in \hat{\kc}_d$ and that $A$ is completely dependent. In fact, the vector 
    $\mathbf{X}$ even fulfills that each coordinate $X_j$ is a function of the other $(d-1)$ coordinates, i.e., a function of $\mathbf{X}_{-j}$. In other words, $A$ is completely 
    dependent in each direction.     
    To complete the proof it suffices to show that all $(d-1)$-dimensional marginals 
    of $A$ coincide with $\Pi_{d-1}$, which can be done as follows. 
    Considering disintegration (w.r.t. the coordinates $2,\ldots,d-1$) and again using  $\lambda^{r_a}=\lambda$,  
    for arbitrary $E_2,E_3,\ldots, E_d \in \mathcal{B}([0,1])$ we get
    \begin{align*}
        \mathbb{P}^{\mathbf{X}_{-1}} \Bigg(\bigtimes_{i=2}^d &E_i\Bigg) = 
        \mathbb{P}^{\mathbf{X}} \left([0,1] \times \bigtimes_{i=2}^d E_i\right) \\
        &= 
        \int_{\bigtimes_{i=2}^{d-1}E_i} \mathbb{P}\left(X_1 \in [0,1], r_{\sum_{i=2}^{d-1}x_i \,\, \textup{mod (1)}}(X_1) \in E_d \right)d\lambda_{d-2}(\mathbf{x}_{2:d-1}) \\
        &= \int_{\bigtimes_{i=2}^{d-1}E_i} \mathbb{P}\left(r_{\sum_{i=2}^{d-1}x_i \,\, \textup{mod (1)}}(X_1) \in E_d \right)d\lambda_{d-2}(\mathbf{x}_{2:d-1}) \\
        &= \int_{\bigtimes_{i=2}^{d-1}E_i} \lambda(E_d) \,d\lambda_{d-2}(\mathbf{x}_{2:d-1}) 
        = \prod_{i=2}^d \lambda(E_i) = \prod_{i=2}^d \mathbb{P}^{X_i}(E_i). 
        \end{align*}
    This shows $\mathbf{X}_{-1} \sim \Pi_{d-1}$. Replacing $j=1$ by any other index
    $j \in \{2,\ldots,d-1\}$ and proceeding analogously yields $\mathbf{X}_{-j} \sim \Pi_{d-1}$. 
    Since $\mathbf{X}_{-d} \sim \Pi_{d-1}$ by construction, the proof is complete. The left panel of Figure \ref{fig1} depicts a sample of the copula $A$ for 
    the case $d=3$. 
    \end{proof}

\end{document}
